\documentclass[12pt,a4paper]{article}
\usepackage[utf8]{inputenc}
\usepackage[T1]{fontenc}
\usepackage[brazil]{babel}
\usepackage{geometry}
\usepackage{setspace}
\usepackage{indentfirst}
\usepackage{cite}
\usepackage{hyperref}

\geometry{
top=3cm,
bottom=2cm,
left=3cm,
right=2cm
}

\onehalfspacing

\title{
\textbf{Sistema de Apoio à Priorização da Fiscalização Ambiental de Imóveis Rurais utilizando Machine Learning}
}

\author{
Bruno Schramm Vendruscolo 
Pedro Henrique Klein
}

\date{}

\begin{document}

\maketitle

\section{Introdução}

O monitoramento ambiental por meio de dados geoespaciais e sensoriamento remoto tem possibilitado a identificação de alterações na cobertura vegetal em grandes extensões territoriais\cite{monitoramentdesmatamento}. No Brasil, iniciativas como o MapBiomas Alertas disponibilizam informações sobre áreas com ocorrências de desmatamento, enquanto o Cadastro Ambiental Rural (CAR) fornece dados georreferenciados sobre imóveis rurais.

A grande quantidade de dados, dificulta a definição de quais áreas devem receber maior atenção das equipes de fiscalização. Com isso o uso de técnicas de Machine Learning podem auxiliar na análise de diferentes características dos imóveis e de seu histórico ambiental, permitido assim identificar padrões e estabelecer critérios de priorização de fiscalização.

Diante disso o trabalho propõe o desenvolvimento de uma ferramenta capaz de integrar dados públicos do CAR e do MapBiomas para auxiliar na priorização de imóveis rurais para análise e fiscalização ambiental.

\section{Problema}

Identificar alertas de desmatamento não classifica qual a ocorrência de maior urgência ou que possui um maior impacto ambiental, com isso torna-se necessário desenvolver mecanismos que auxiliem na classificação de fiscalizações ambientais.

O Instituto Brasileiro do Meio Ambiente e dos Recursos Naturais Renováveis (IBAMA) estabelece a utilização de análise de riscos no planejamento das ações de fiscalização ambiental, indicando a necessidade de estabelecer prioridades entre as ocorrências. \cite{ibamalegislacao}

\section{Objetivos}

\subsection{Objetivos Gerais}
Desenvolver uma ferramenta para auxiliar na classificação de prioridade de alertas de desmatamento, utilizando Machine Learning.

\subsection{Objetivos Específicos}

\begin{itemize}
    \item Ajudar na classificação de prioridades de fiscalizações de crimes ambientais;
    \item Relacionar dados do SiCAR com Alertas disponíveis no MapBiomas Alertas;
    \item Desenvolver um modelo capaz de priorizar alertas de desmatamento;
    \item Desenvolver uma aplicação em Streamlit suportando o uso do modelo.
\end{itemize}

\section{Justificativa}

A fiscalização ambiental desempenha um papel importante no controle do desmatamento. Um estudo recente mostra que fiscalizar e aplicar as leis na legislação ambiental podem contribuir para a redução do desmatamento na Amazônia brasileira. \cite{amazoniabrasileira}

Diante disso, a capacidade de fiscalizar é limitada diante da extensão territorial e da quantidade de ocorrencias identificadas. O IBAMA reconhece essa necessidade ao estabelecer a análise de riscos e a definição de prioridades como parte o planejamento de ações de fiscalização. \cite{ibamalegislacao}

Outro estudo relacionado a monitoramento de desmatamento destaca o desafio envolvendo o processamento e analise de grandes volumes de dados, vindo de sensoriamento remoto. \cite{monitoramentdesmatamento}

Com isso técnicas de Machine Learning podem auxiliar na identificação de padrões presentes nos dados e nos imóveis rurais classificando assim uma prioridade de fiscalização. Portanto o trabalho busca ajudar a transformar dados ambientais públicos em informações que possam ajudar na fiscalização do combate ao desmatamento.

\begin{thebibliography}{9}

\bibitem{monitoramentdesmatamento}
\textit{Enhancing deforestation monitoring in the Brazilian Amazon: A semi-automatic approach leveraging uncertainty estimation}.
2024.
Disponível em:
\url{https://www.sciencedirect.com/science/article/abs/pii/S0924271624000765}.

\bibitem{ibamalegislacao}
IBAMA.
\textit{Portaria nº 217, de 10 de outubro de 2023}.
Instituto Brasileiro do Meio Ambiente e dos Recursos Naturais Renováveis, 2023.
Disponível em:
\url{https://www.ibama.gov.br/component/legislacao/?legislacao=139359\&view=legislacao}.

\bibitem{amazoniabrasileira}
\textit{Environmental enforcement and sustainable development: Causal evidence at the municipal level for the Brazilian Amazon}.
Journal of Environmental Management, 2026.
Disponível em:
\url{https://www.sciencedirect.com/science/article/abs/pii/S0301479726006468}.

\end{thebibliography}

\end{document}